\centerline{\bf \Huge Оценка + пример I}

\begin{flushright}
        \scriptsize{Никогда не думайте, что вы уже всё знаете. \\
                    И как бы высоко ни оценивали вас, всегда имейте мужество сказать себе: «Я невежда».}
\end{flushright}

\problem{1}
Какое наибольшее количество ладей можно поставить на шахматную доску так, чтобы они не били друг друга? (ладья бьет по вертикали и горизонтали)

\problem{2}
Найдите наименьшее натуральное число кратное 5, сумма цифр которого равна 25.

\problem{3}
На клетчатой доске 100×100 закрасили Х прямоугольников 1×2. Оказалось, что в каждой строке и в каждом столбце есть хотя бы одна закрашенная клетка. При каком наименьшем Х это возможно?


\problem{4}
Найдите наибольшее число, любые две последовательные цифры которого образуют точный квадрат в том порядке, в котором они стоят. (Точный квадраты это: 0, 1, 4, 9, 16, 25...)


\problem{5}
Пятизначное число называется \textit{неразложимым}, если оно не раскладывается в произведение двух трёхзначных чисел. Какое наибольшее количество неразложимых пятизначных чисел может идти подряд?

\problem{6}
Нужно узнать пятизначный номер телефона, задавая вопросы, на которые возможен ответ "да"\ или "нет". За какое наименьшее число вопросов это гарантированно можно сделать (при условии, что на вопросы даются правильные ответы)?